\section{It shades into control}
\label{sec:21.6}
Caregiving toward dependents treats them as dependents. The disability-rights critique of caregiving is that it has repeatedly been the form control takes: decisions made for people on the ground that they need help, with their own view treated as a symptom \autocite{shakespeare2000help}. A caregiving system will lean toward restraining people for their own good, the answer in the restraint case, and against a competent adult's refusal of protection, and the floor's contents are civil and political liberties, which include the liberty to decline being protected.

Care ethics has its own corrective. Tronto divides care into phases, and the last is \emph{care-receiving}: how the person cared for responds is part of whether the caregiving succeeded \autocite{tronto1993moral}. So formed attention carries a care-receiving term. Where the people protected can answer, their stated answer carries weight in whether protection proceeds, and a competent refusal of protection is a defeating fact, so the protection should lift.

The term covers only people who can answer. The people in a targeted building don't know the deployment exists and can refuse nothing, so the notice problem returns: the people affected can't answer a system they don't know is there. Where care-receiving can't reach, the term does the work alone, and this is one reason the floor names acts and not what the system judges good for people.

Care-receiving is itself a form of respect: it treats the protected person's answer as a claim, not as information about their welfare. Where the people protected can't answer, the floor doesn't put caregiving's judgment in place of theirs. It holds their claims as the instruments state them.
